% answers_short.tex —— 1994 年·上海试点教材理科卷 · 简短答案速查
% 由 tools/make_short_answers.py 自 answers.tex 的 \gksolhead/\gkanswer 提取，
% 并与 1994/1994_数学真题_上海试点理_答案要点.md 逐题比对后生成，勿手改（改源 MD 后重跑）。

\begin{gknotice}
\textbf{使用说明}：本节只给最终答案便于快速核对，完整推导见解析版。答案据固定来源 PDF 后附解答及评分标准逐题转录（出版者、版次及答案官方性未核实）。\textbf{非官方标准答案}。
\end{gknotice}

\gkpart{填空题}
\begin{itemize}[leftmargin=2.6em]
  \item[\textbf{1.}] \(-2<x<0\)
  \item[\textbf{2.}] \(-4\)
  \item[\textbf{3.}] \((x+2)^2+(y-3)^2=4\)
  \item[\textbf{4.}] \(x=4\)
  \item[\textbf{5.}] \(b=-\dfrac23\)
  \item[\textbf{6.}] \(y=-\sqrt{x^2+1}\ (x\ge0)\)
  \item[\textbf{7.}] \((0,\sqrt3)\) 和 \((0,-\sqrt3)\)
  \item[\textbf{8.}] \(40\)
  \item[\textbf{9.}] \(300\)
  \item[\textbf{10.}] \(\dfrac\pi2\)
\end{itemize}

\gkpart{选择题}
\begin{center}\begin{tabular}{ccccccccccc}
\toprule
题号 & 11 & 12 & 13 & 14 & 15 & 16 & 17 & 18 & 19 & 20 \\
\midrule
答案 & B & A & C & C & D & B & A & D & B & C \\
\bottomrule
\end{tabular}\end{center}

\gkpart{解答题}
\begin{itemize}[leftmargin=2.6em]
  \item[\textbf{21.}] \(\tan(\alpha-2\beta)=\dfrac7{24}\)
  \item[\textbf{22.}] (1) \(\theta=\arccos\left(-\dfrac{\sqrt{10}}{10}\right)\)；(2) \(k=-\dfrac52\) 或 \(k=2\)
  \item[\textbf{23.}] (1) \(\arctan\dfrac{\sqrt5}{3}\)；(2) \(\dfrac{\sqrt2}{2}a\)
  \item[\textbf{24.}] (1) \(\dfrac23t^3\)；(2) \(P=\left(-\dfrac t3,\dfrac{t^2}3\right)\)，\(M=\left(\dfrac t3,\dfrac53t^2\right)\)，轨迹为 \(y=15x^2\)
  \item[\textbf{25.}] (1) $S(t)=\begin{cases}2(1-t+t^2-t^3),&0<t<\dfrac12,\\[2pt]\dfrac12\left(t+\dfrac1t\right),&t\ge\dfrac12.\end{cases}$；(2) 递减区间 \((0,1]\)，递增区间 \([1,+\infty)\)
  \item[\textbf{26.}] (1) \(0<q<\dfrac{1+\sqrt5}{2}\)；(2) $b_n=(1+r)q^{n-1}$，$\displaystyle\lim_{n\to\infty}\frac1{S_n}=\begin{cases}\dfrac{1-q}{1+r},&0<q<1,\\[2pt]0,&q\ge1.\end{cases}$；(3) 最大项为 \(2.25\)，最小项为 \(-4\)
\end{itemize}
